\section{Additional source results}\label{app:additional-source-results}
\subsection{Effect of the latent spring source}
With backbone, initialization, and data fixed, harmonic improves graph validity over isotropic Gaussian by 5.86 points on 12 development compositions, 8.83 on 10 additional compositions, and 8.20 on the 24-composition panel. It also outperforms the tested covariance-matched Gaussian approximation. Appendix~\ref{app:source-panels} gives the Gaussian, shell-tree, and covariance controls.

\subsection{Generalization across molecular sizes}\label{sec:wide}
We evaluate the frozen source models on 64 further compositions, with 16 in each
of four size ranges from 17 to 64 atoms. All are absent from both processed OMol25
corpora and from the earlier generation panels. The models receive no additional
training. Metadata ranks select each composition after the same reference-validity
checks, with 32 draws per composition and model continuation.

\begin{table}[t]
\centering
\caption{Raw chemical-graph validity on the 64-composition panel. Each size range
contains 1,024 attempts per method across two continuations. Values are percentages;
the final column reports percentage-point differences.}
\begin{tabular}{lrrr}
\toprule
Atoms & Gaussian & Harmonic & Difference\\
\midrule
17--28 & 49.12 & 55.08 & +5.96\\
29--40 & 24.02 & 38.28 & +14.26\\
41--52 & 16.50 & 21.29 & +4.79\\
53--64 & 9.77 & 12.40 & +2.64\\
\midrule
All 64 compositions &24.85&31.76&+6.91\\
\bottomrule
\end{tabular}\label{tab:wide}
\end{table}

The harmonic source improves the pooled rate in every size range. Overall it
adds 6.91 percentage points, with a size-stratified composition-bootstrap 95\%
interval of [4.74,9.03]. The two continuations improve by 11.18 and 2.64 points,
respectively. Absolute validity decreases for larger molecules, while the source
advantage remains positive in the aggregate. These results extend the source
comparison beyond the original 17--28-atom confirmation range.

\subsection{Independent backbone and generator objectives}\label{sec:matched-generators}
Gaussian FM, harmonic FM, position-conditional EDM, and GAGA \citep{edm,qu2026gaga} train from scratch on the same EGNN and 20,000 OMol25 structures. Two repetitions share initialization within each comparison, with 30,000 updates at batch size 32. Neither self-conditioning nor physical adaptation is used. The FM models differ only in source, sharing interpolation and atom pairing (Appendix~\ref{app:matched-generators}).

Harmonic improves FM graph validity by 1.34 points (composition 95\% interval [0.12,2.52]). EDM and GAGA achieve higher validity in this setting, with different paths, objectives, and samplers. Distance self-conditioning subsequently supplies the FM parent for physical-head comparisons at matched training-forward budgets (Appendices~\ref{app:feedback-transfer} and~\ref{app:matched-connection}).

A separate 16-composition complete-system comparison gives 38.87\% GFN2 joint yield for the inherited 5.90-million-parameter FlowMol backbone and 46.88\% with force correction. The 2.38-million-parameter GAGA trained from scratch gives 9.77\% at 128 calls and 12.11\% at 651. Appendix~\ref{app:common-models} details the differing training histories and physical supervision, costs, capacity controls, and fitting-duration study.


\paragraph{Physical weight updates.}
On 24 new compositions, paired physical training lowers energy and force RMS under both eSEN and independent GFN2 evaluation, with similar graph validity (Appendix~\ref{sec:fresh-physics}). Both fits also favor the paired update over a norm-matched direct update (Appendix~\ref{sec:update-ablation}). Appendix~\ref{sec:absolute-quality} compares physical adaptation under both sources.

\paragraph{Training components.}
At matched training-forward budgets, self-conditioning improves validity by 12.43 points ($[9.70,15.23]$). Removing the force displacement raises energy on paired graph-valid outputs by 0.01939 eV/atom. Appendix~\ref{app:components} reports both runs and the update-matched controls.
